\documentclass{article}

\usepackage[a4paper, left=1.5cm, right=1.5cm, top=2cm, bottom=2cm]{geometry}

\usepackage{../../../../components/components}
\usepackage{hyperref}
\usepackage{fancyhdr}


% Configuration des en-têtes et pieds de page
\pagestyle{fancy}
\fancyhf{} % reset tout

\fancyhead[L]{DL2 Math-Info ASF4}
\fancyhead[C]{Algèbre}
\fancyhead[R]{2025-2026}

\fancyfoot[L]{Ewen Rodrigues de Oliveira}
\fancyfoot[R]{\thepage}

\begin{document}

\docTitle{Chapitre 10 : Formes bilinéaires et formes quadratiques}

\warning{
    Cours non finalisé, non structuré et contenant beaucoup d'oublis.\\
    Il est pour l'instant peu sûr de réviser à partir de ce cours, mais il peut être intéressant de le parcourir pour se faire une idée du sujet.
}
\section{Formes bilinéaires}
\subsection{Généralités}
\definition{
    Soit $E$ un $\mathbb{R}$-espace vectoriel.\\
    Une \textbf{forme bilinéaire} sur $E$ est une application $f : E \times E \to \mathbb{R}$ telle que :
    \begin{itemize}
        \item pour tout $x \in E$, l'application $y \mapsto f(x,y)$ est linéaire,
        \item pour tout $y \in E$, l'application $x \mapsto f(x,y)$ est linéaire.
        \item autrement dit, $f$ est linéaire par rapport à chacune de ses variables.
    \end{itemize}
}
\vocabulary{On note $\mathcal{B}(E)$ l'ensemble des formes bilinéaires sur $E$.}
\example{Le produit scalaire est une forme bilinéaire.}

\subsection{Formes symétriques}

\definition{
    On dit qu'une forme bilinéaire $f$ est \textbf{symétrique} si $f(x,y) = f(y,x)$ pour tous $x,y \in E$.
}

\definition{
    Soit $\varphi : E \times E \to \mathbb{R}$.\\
    L'application $D_\varphi : E \to E^*$ telle que $\forall x,y \in E, \; \varphi(x,y) = \langle D_\varphi(x), y \rangle$ est appelée l'\textbf{application linéaire canonique} de $\varphi$.
}

\subsection{Changement de base}
\theorem{Proposition}{Changement de base}{false}{
    Soient $B=(e_1,\ldots,e_n)$ et $B'=(e_1',\ldots,e_n')$ deux bases de $E$, et $P_{B\to B'}$ la matrice de passage de $B$ à $B'$.\\
    Soit $f$ une forme bilinéaire. Alors \[ Mat_{B'}(f) = {}^t P_{B\to B'} Mat_B(f) P_{B\to B'} \]
}

\noindent{\textbf{Démonstration :}\\
    Soient $x, y\in E$, on a $P_{B\to B'}$ la matrice de passage de $B$ à $B'$. \\
    Posons $X$ la matrice de $x$ dans la base $B$ ; $X'$ la matrice de $x$ dans la base $B'$ ;\\
    $Y$ la matrice de $y$ dans la base $B$ ; $Y'$ la matrice de $y$ dans la base $B'$. \\
    La matrice de passage $P$ satisfait $ X=PX'$ et $Y=PX'$.\\
    On sait de plus que $f(x, y) = {}^t XMat_B(f)Y = {}^t X' Mat_{B'}(f)Y'$.\\
    Donc ${}^tX Mat_B(f)Y = {}^t X' Mat_{B'}(f) Y'\Leftrightarrow {}^tP\cdot Mat_B(f)\cdot PY'={}^tX'\cdot Mat_{B'}(f)Y'$.\\
    Donc, ${}^tP\cdot Mat_B(f)\cdot P = Mat_{B'}(f)$.
}

\subsection{Rang et noyau d'une forme bilinéaire symétrique}

\definition{
    Soit $f$ une forme bilinéaire symétrique sur $E$, alors : 
    \begin{enumerate}
        \item $ker(f) = \{ x \in E \mid \forall y \in E f(x,y) = 0 \}$,
        \item $rg(f) = \dim(E) - \dim(ker(f))$.
    \end{enumerate}
}
\vocabulary{On dit que $f$ est \textbf{non-dégénérée} si $ker(f) = \{0\}$, c'est à dire si $rg(f) = \dim(E)$.}

\theorem{Théorème}{}{false}{
Soit $(e_1, \ldots, e_n)$ une base de $E$, et $A_{ij} = f(e_i, e_j)$, alors on a :
\begin{enumerate}
    \item $f$ est non dégénérée si et seulement si $A$ est inversible,
    \item $ker(f) = \{ x \in E \mid A X = 0 \}$,
    \item $rg(f) = rg(A)$.
\end{enumerate}
}

\noindent{\textbf{Démonstration :}\\
Posons $f : E \times E \to \mathbb{R}$ telle que $f : x \mapsto (y \mapsto f(x,y))$.
Posons aussi $d : E \to E*$.\\
Comme $f$ est bilinéaire, $d$ est linéaire.\\
Ainsi, $ker(f) = ker(d)$ et $rg(f) = rg(d)$.\\
Comme $A$ est la matrice de $d$ dans les bases $(e_1, \ldots, e_n)$ et $(e_1^*, \ldots, e_n^*)$, on a $ker(d) = \{ x \in E \mid A X = 0 \}$ et $rg(d) = rg(A)$.
}

\example{
    Un produit scalaire est une forme bilinéaire symétrique non dégénérée.\\
    En effet, si $f$ est un produit scalaire, alors $f(x,y) = f(y,x)$ pour tous $x,y \in E$, et $ker(f) = \{0\}$ (car $f(x,x) > 0$ pour tout $x \neq 0$).
}


\theorem{Proposition}{}{false}{
    Soit $f$ une forme bilinéaire symétrique. Alors :
    \begin{enumerate}
        \item $f$ est non-dégénérée si et seulement si $rg(f) = \dim(E)$,
        \item Soit $B=(e_1,\ldots,e_n)$ une base de $E$. Alors $rg(Mat_B(f)) = rg(f)$.\\
        On a de plus que $f$ est non-dégénérée si et seulement si $Mat_B(f)$ est inversible.
    \end{enumerate}
}

\subsection{Orthogonalité}


\definition{
    Soit $\varphi$ une forme bilinéaire symétrique sur $E$ un $\mathbb{R}$-espace vectoriel.
    \begin{enumerate}
        \item Deux vecteurs $x, y \in E$ sont \textbf{orthogonaux} si $\varphi(x,y) = 0$, on note $x \perp_\varphi y$.
        \item L'orthogonal d'une partie $H \subseteq E$ est le sous-espace vectoriel $H^\perp = \{ z \in E \mid \forall x \in H, x \perp_\varphi z \}$.\\
        Si $\varphi$ est un produit scalaire, alors $H^\perp \cap H = \{0\}$.
        \item Un vecteur $x \in E$ est dit \textbf{isotrope} si $x \perp_\varphi x$.
    \end{enumerate}
}

\remark{Une forme bilinéaire symétrique qui n'a pas de vecteur isotrope est nécessairement non dégénérée.}

\theorem{Proposition}{}{false}{
    Soit $f$ une forme bilinéaire symétrique. Soit $F$ un sous-espace vectoriel de $E$.\\
    Alors :
    \[
        dim(F^\perp_f) = dim(E) - dim(F) + dim(F \cap ker(f))
    \]
    Si $f$ est non dégénérée, alors $dim(F^\perp_f) = dim(E) - dim(F)$.
}

\noindent{\textbf{Démonstration :}\\
    On regarde l'application $D : \begin{matrix}    E\to E^*\\
                                                    x\mapsto (y\mapsto f(x, y)) \end{matrix}$.\\
    $D$ est une application linéaire car $f$ est bilinéaire. \\
    $x\in Ker(D)$, donc $D(x)=0$ si et seulement si $\forall y\in E, f(x, y)=0$. Donc $Ker(D)=Ker(f)$. \\
    Si $B$ est non-dégénérée, alors $Ker(D)={0}$. Donc $D$ est un isomorphisme.\\
    Soit $L={\ell\in E^* : \ell_F|=0}$.\\
    Donc, $D(F^\bot)=L$. Donc $dim(F^{\bot_f})=dim(L)=dim(E)-dim(F)$. (d'après la remarque en dessous)\\
}

\remark{ $D(x)\in L\Leftrightarrow \forall y\in F, D(x)(y)=0\Leftrightarrow f(x, y)=0 \forall y\in F \Leftrightarrow x\in F^\bot$.}	

\theorem{Corollaire}{}{false}{
    Si $f$ est une forme bilinéaire symétrique non dégénérée sur $E \times E$, avec $E$ de dimension finie, alors $(F^\perp_f)^\perp = F$ pour tout sous-espace vectoriel $F$ de $E$.
}

\noindent{\textbf{Démonstration : }\\
    $dim((F^{\bot_f})^{\bot_f})=dim(E)-dim(F^{\bot_f})=dim(E)-(dim(E)-dim(F))=dim(F)$.\\
    L'inclusion $F\subseteq (F^{\bot_f})^{\bot_F}$ est automatique car $x\in F$, alors $\forall y\in F^{\bot_f}, f(x, y)=0$, alors $x\in (F^{\bot_f})^{\bot_f}$.\\
}

\subsection{Restriction d'une forme bilinéaire symétrique}
Si $b$ est un produit scalaire sur $E$, il est clair que la restriction de $b$ à tout sous-espace vectoriel $F$ de $E$ est encore un produit scalaire sur $F$. En revanche, l'exemple suivant montre qu'une forme bilinéaire symétrique non dégénérée peut avoir des restrictions qui sont
dégénérées.

\example{
    Soit $E = \mathbb{R}^2$ et $b$ la forme bilinéaire symétrique définie par $b((x_1, y_1), (x_2, y_2)) = x_1 x_2 - y_1 y_2$.\\
    Alors $b$ est non dégénérée, mais la restriction de $b$ à la droite vectorielle $F = \{ (x, x) \mid x \in \mathbb{R} \}$ est dégénérée.
}

\theorem{Proposition}{}{false}{
    Soit $E$ un $\mathbb{R}$-espace vectoriel de dimension finie, et soit $\varphi$ une forme bilinéaire symétrique sur $E \times E$.\\
    Soit $F$ un sous-espace vectoriel de $E$. Alors :
    \[
         \varphi_{|F \times F} \text{ est non dégénérée } \Leftrightarrow F \cap F^\perp_f = \{0\} \Leftrightarrow F \oplus F^\perp_f = E
    \]
}
\noindent{\textbf{Démonstration : }
    \begin{itemize}
        \item "$\varphi_{|F\times F}$ est non dégénérée $\Rightarrow F\cap F^{\bot_f} = \{0\}$" : Soit $x\in F\cap F^\bot$, alors $\forall y\in F, \varphi(x, y)=0$\\
        Donc $x\in Ker(f_{|F\times F})$. Donc, si $f_{|F\times F}$ est non-dégénérée, alors $x=y$.
        \item "$F\cap F^{\bot_f}={0}\Rightarrow F\oplus F^{\bot_f}=E$ " : $dim(F^\bot)=dim(E)-dim(F)+dim(F\cap Ker(f))$. \\
        Si $x\in F\cap Ker(f)$, alors $f(x, y)=0, \forall y\in E$. Donc, $x\in F^\bot$, donc $x\in F\cap F^\bot=\{0\}$, donc $x=0$.\\
        Donc, $dim(F\cap Ker(f))=0$. \\
        On a donc, $dim(F^\bot)=dim(E)-dim(F)$. Donc $dim(E)=dim(F)+dim(F^\bot)$ et $F\cap F^\bot=\{0\}$. \\
        Ainsi, $E= F\oplus F^\bot$.
        \item "$F\oplus F^\bot\Rightarrow \varphi_{|F\times F}$" : Si $x\in Ker(\varphi_{|F\times F})$. Donc, $x\in F$ et $\varphi(x, y)=0, \forall y\in F$. Donc, $x\in F$ et $x\in F^(\bot_\varphi)$. \\
        Donc $x\in F\cap F^{\bot_\varphi}$. Mais si $E=F\cap F^{\bot_\varphi}$, alors $F\cap F^{\bot_\varphi}=0$.
    \end{itemize}

}

\section{Formes quadratiques}
Dans le cadre des produits scalaires, on a vu que la fonction "carré de la norme" avait des propriétés remarquables, en particulier vis-à-vis de l'orthogonalité grâce au théorème de Pythagore. La généralisation de cette fonction dans le cadre des formes bilinéaires conduit à la notion de forme quadratique.

\subsection{Généralités}

\definition{
    Soit $E$ un $\mathbb{R}$-espace vectoriel.\\
    Une fonction $q : E \to \mathbb{R}$ est une \textbf{forme quadratique} sur $E$ s'il existe une forme bilinéaire symétrique $\varphi$ sur $E \times E$ \textit{(a priori non symétrique)} telle que $q(x) = \varphi(x,x)$ pour tout $x \in E$.
}

\remark{Toute forme quadratique $q$ vérifie : $q(\lambda x) = \lambda^2 q(x)$ pour tout $\lambda \in \mathbb{R}$ et tout $x \in E$. En particulier, $q(0) = 0$.}

% Retour de Romain Petrides.

\example{
    Dans $\mathbb{R}^4$, $q : (x,y,z,t) \mapsto x^2 + y^2 + z^2 - x^2 \cdot t^2$ est une forme quadratique.\\
    En effet, on peut définir $\varphi((x_1, y_1, z_1, t_1), (x_2, y_2, z_2, t_2)) = x_1 x_2 + y_1 y_2 + z_1 z_2 - c^2t_1t_2$.\\
    On a que $q(x,y,z,t) = \varphi((x,y,z,t), (x,y,z,t))$ pour tout $(x,y,z,t) \in \mathbb{R}^4$.
}

\subsection{Polarisation d'une forme quadratique}

\theorem{Proposition}{}{false}{
    Soit $q : E \to \mathbb{R}$ une forme quadratique.\\
    Alors, $\exists ! \varphi \in \mathcal{B}(E)$ symétrique telle que $\forall  x\in E, \; q(x) = \varphi(x,x)$.
}

\noindent{  \textbf{Démonstration : }
    \begin{itemize}
        \item \underbar{Analyse} : Soit $\varphi$ une forme bilinéaire symétrique telle que $q(x)=\varphi(x, x)$.\\
        Alors, $q(x+y) = \varphi(x+y, x+y) = \varphi(x, x)+\varphi(x, y)+\varphi(y, x)+\varphi(y, y) = \varphi(x, x)+2\varphi(x, y)+\varphi(y, y)$\\
        $=q(x)+2\varphi(x, y)+q(y)$. Donc, $\varphi(x, y)=\frac{q(x+y)-q(x)-q(y)}{2}$.\\
        Si une telle forme bilinéaire symétrique existe, alors elle est unique.
        \item \underbar{Synthèse} : On pose $\varphi(x, y)=\frac{q(x+y)-q(x)-q(y)}{2}$. Il reste à montrer que $\varphi$ est bilinéaire.\\
        Soit $\psi\in\mathcal{B}(E)$ telle que $\forall x\in E, q(x)=\psi(x, x)$.\\
        On a $\varphi(x, y)=\frac{q(x+y)-q(x)-q(y)}{2}=\frac{\psi(x+y, x+y)-\psi(x, x)-\psi(y, y)}{2}$ et $\varphi\in\mathcal{B}(E)$ car $\varphi$ est combinaison linéaire de $\psi(x, y)$ et $\psi(y, x)$.
    \end{itemize}
}

\remark{$\mathcal{B}(E)$ désigne l'ensemble des formes bilinéaires sur $E$.}
\vocabulary{On appelle $\varphi$ la \textbf{forme polaire} de $q$.}

\example{
    Soit $(x,y,z) \in \mathbb{R}^3$.\\
    Considérons $q(x,y,z) = xy + 2x^2 - y^2 + 3zx - zy$.\\
    La forme polaire de $q$ est $\varphi : ((x_1, y_1, z_1), (x_2, y_2, z_2)) \mapsto 2x_1x_2 - y_1 y_2 + \frac{x_1y_2+y_1x_2}{2} + 3 \cdot \frac{z_1x_2+x_1z_2}{2} - \frac{z_1y_2+z_2y_1}{2}$.\\
    De plus, si $e$ est la base canonique de $\mathbb{R}^3$, alors $Mat_e(q) = Mat_e(\varphi) = \begin{pmatrix}
2 & \frac{1}{2} & \frac{3}{2} \\
\frac{1}{2} & -1 & -\frac{1}{2} \\
\frac{3}{2} & -\frac{1}{2} & 0
\end{pmatrix}$. 
}

\definition{
    Soit $q : E \to \mathbb{R}$ une forme quadratique sur $E$, où $E$ est un $\mathbb{R}$-espace vectoriel de dimension finie muni d'une base $e$.\\
    \begin{itemize}
        \item $Mat_e(q) = Mat_e(\varphi) = Mat_(e,e^*)(D_\varphi)$, où $\varphi$ est la forme polaire de $q$.
        \item $Ker(q) = Ker(\varphi) = Ker(D_\varphi)$,
        \item $rg(q) = rg(\varphi) = rg(D_\varphi)$.
        \item $q$ est \textbf{dégénérée} si $Ker(D_\varphi) = \{0\}$.
    \end{itemize}
}

\ndlr{Vérifier le dernier point, semble être une erreur de frappe, devrait être "non dégénérée" au lieu de "dégénérée".}
\ndlr{Qu'est-ce que $D_\varphi$ ?}

\remark{On observe que $Mat_e(q)$ est symétrique.}

\subsection{Bases q-orthogonales}
\subsubsection{Définitions et propriétés}
\definition{
    Soient $x,y \in E$ et $q$ une forme quadratique sur $E$.\\
    \begin{itemize}
        \item $x$ est \textbf{isotrope} si $q(x) = 0$,
        \item On dit que $x$ et $y$ sont \textbf{orthogonaux} si $\varphi(x,y) = 0$, où $\varphi$ est la forme polaire de $q$.
    \end{itemize}
}
\remark{$x$ est isotrope si et seulement si $x \perp x$.}

\ndlr{Revoir la définition ci-dessus, semble être un duplicata de ce qu'on a vu auparavant.}

\definition{
    Soit $E$ un $\mathbb{R}$-espace vectoriel de dimension finie, et $q$ une forme quadratique sur $E$.\\
    Considérons $e=(e_1, \ldots, e_n)$ une base de $E$.\\
    On dit que $e$ est une \textbf{base $q$-orthogonale} si les éléments de $e$ sont deux à deux orthogonaux, c'est à dire si $\varphi(e_i, e_j) = 0$ pour tous $i \neq j$, où $\varphi$ est la forme polaire de $q$.
}

\remark{
    Si $e$ est une base $q$-orthogonale de $E$, alors $Mat_e(q)$ est une matrice diagonale.
}

\remark{
    L'application $q = N^2$ où $N$ est une norme sur $E$ est une forme quadratique.
}

\example{
    Considérns $q(x,y,z) = x^2 - y^2 + 2z^2$ une forme quadratique sur $\mathbb{R}^3$.\\
    Alors, la base canonique de $\mathbb{R}^3$ est une base $q$-orthogonale. En effet, les éléments de la base canonique sont deux à deux orthogonaux, et on a $q(e_1) = 1$, $q(e_2) = -1$ et $q(e_3) = 2$.
}

\subsubsection{Détermination d'une base $q$-orthogonale}

\training{
Si $q(x) = e_1(\ell_1(x))^2 + e_2(\ell_2(x))^2 + \cdots + e_k(\ell_k(x))^2$ où $(\ell_1, \ldots, \ell_k)$ une famille libre de $E^*$ et $c_1, \ldots, c_k \in \mathbb{R}$.

\begin{enumerate}
    \item On complète $(\ell_1, \ldots, \ell_k)$ en une base de $E^*$, disons $(\ell_1, \ldots, \ell_n)$.
    \item On calcule la base antéduale de $(\ell_1, \ldots, \ell_n)$, disons $(e_1, \ldots, e_n)$.
\end{enumerate}
Alors, $(e_1, \ldots, e_n)$ est $q$-orthogonale.\\
On a $\langle \ell_j, e_i \rangle = \delta_{ij}$. Soit $\varphi$ la forme linéaire de $q$.\\
Alors $\varphi(x,y) = \sum_{s=1}^k c_s \ell_s(x) \ell_s(y)$ et $\varphi(e_i, e_j) = \sum_{s=1}^k c_s \ell_s(e_i) \ell_s(e_j) = \sum_{s=1}^k c_s \delta_{si} \delta_{sj} = 0$ si $i \neq j$.
}
\example{
    Soit $q(x,y,z) = x^2 + (y+z)^2$ une forme quadratique sur $\mathbb{R}^3$.\\
    On a :
    \[
        \ell_1(x,y,z) = x, \quad \ell_2(x,y,z) = y+z, \quad \ell_3(x,y,z) = z.
    \]
    On a de plus que $P_{can \to (\ell_1, \ell_2, \ell_3)} = \begin{pmatrix}
1 & 0 & 0 \\
0 & 1 & 0 \\
0 & 1 & 1
\end{pmatrix}$, et par pivot de Gauss, on trouve que $P_{can \to (\ell_1, \ell_2, \ell_3)}^{-1} = \begin{pmatrix}
1 & 0 & 0 \\    
0 & 1 & -1 \\
0 & 0 & 1
\end{pmatrix}$.\\
    Donc, la base antéduale de $(\ell_1, \ell_2, \ell_3)$ est $(e_1, e_2, e_3 - e_2)$, et cette base est $q$-orthogonale.
}

\theorem{Proposition}{}{false}{
    Soit $E$ un $\mathbb{R}$-espace vectoriel de dimension finie, et $q$ une forme quadratique sur $E$.\\
    Considérons $q(x) = \sum_{i=1}^k c_i (\ell_i(x))^2$ où $(\ell_1, \ldots, \ell_k)$ une famille libre de $E^*$ et $c_1, \ldots, c_k \in \mathbb{R}$.\\ 
    Alors :
    \begin{enumerate}
        \item La forme quadratique $q$ est de rang $k$,
        \item Soit $(l_{k+1}, \ldots, l_n)$ une famille de $E^*$ telle que $(\ell_1, \ldots, \ell_k, l_{k+1}, \ldots, l_n)$ soit une base de $E^*$. On note $e$ la base antéduale de $(\ell_1, \ldots, \ell_k, l_{k+1}, \ldots, l_n)$. Alors $e$ est une base $q$-orthogonale de $E$.\\
        De manière équivalente, \[
\mathrm{Mat}_e(q) =
\begin{pmatrix}
c_1 & 0 & \cdots & 0 & 0 & \cdots & 0 \\
0 & c_2 & \cdots & 0 & 0 & \cdots & 0 \\
\vdots & \vdots & \ddots & \vdots & \vdots &  & \vdots \\
0 & 0 & \cdots & c_k & 0 & \cdots & 0 \\
0 & 0 & \cdots & 0 & 0 & \cdots & 0 \\
\vdots & \vdots &  & \vdots & \vdots & \ddots & \vdots \\
0 & 0 & \cdots & 0 & 0 & \cdots & 0
\end{pmatrix}
\in \mathcal{M}_n(\mathbb{R})
\]
    \end{enumerate}
}

\remark{$q$ est une forme quadratique de forme polaire $\varphi(x, y)=\sum_{i=1}^kc_i\ell_i(x)\ell_i(y)$.}

\noindent{\textbf{Démonstration : }\\
    Il suffit de montrer 2. : $2)\Rightarrow rg(q)=rg(...) = k$ ("..." représente la matrice de la proposition)\\
    Montrons 2) : $Mat(q, e)=\sum_{i=1}^kc_s\ell_s(e_i)\ell_(e_j)=\sum_{s=1}^kc_s\delta_{si}\delta_{sj}$. Or, $\delta_{si}\delta_{sj}=1$ si $s=i=j$ et 0 sinon.\\
    Donc, $\varphi(e_i, e_j)=0$ si $i\ne j$ ou $i>k$ ou $j>k$ ; 0 sinon. 
}


\subsubsection{Décomposition de Gauss}
\noindent\href{https://www.youtube.com/watch?v=X9xQ9xiFC9M&t=1s}{Cliquer ici pour une vidéo sur la décomposition de Gauss}

\remark{Le but est d'écrire une forme quadratique quelconque sur $E$ de dimension finie sous la forme $q(x) = \sum_{i=1}^k c_i (\ell_i(x))^2$ où $(\ell_1, \ldots, \ell_k)$ une famille libre de $E^*$ et $c_1, \ldots, c_k \in \mathbb{R}$.}

\example{
    La forme quadratique "contient un carré". On considère $q(x,y,z) = x^2 - 2y^2 + xy - xz$.\\
    On a $q(x,y,z) = (x + \frac{1}{2}y - \frac{1}{2}z)^2 -(\frac{1}{2}y - \frac{1}{2}z)^2 - 2y^2$.\\
    Ici on peut directement poser $\ell_1(x,y,z) = x + \frac{1}{2}y - \frac{1}{2}z$, $\ell_2(x,y,z) = \frac{1}{2}y - \frac{1}{2}z$ et $\ell_3(x,y,z) = y$ pour trouver la base $q$-orthogonale associée à $q$.
}

\example{
    On prend $q(x,y,z) = x^2 - 2xy + xz - y^2 + 2yz$.\\
    On a $q(x,y,z) = (x - y + \frac{1}{2}z)^2 - (-y + \frac{1}{2}z)^2 - y^2 + 2yz - z^2$.\\
    $= (x-y + \frac{1}{2}z)^2 - 2y^2 + 3yz - \frac{5}{4}z^2$.\\
    $= (x-y + \frac{1}{2}z)^2 - 2(y - \frac{3}{4}z)^2 + \frac{9}{8}z^2$ - $\frac{5}{4}z^2$.\\
    $= (x-y + \frac{1}{2}z)^2 - 2(y - \frac{3}{4}z)^2 - \frac{1}{8}z^2$.
}

\example{
    À un moment dans l'algorithme, la forme quadratique "ne contient pas/plus de carré".\\
    Exemple le plus simple : $q(x,y) = xy = \frac{1}{4}((x+y)^2 - (x-y)^2)$.
}

\example{
    $q=(x_1, x_2, x_3, x_4)=x_1x_2+x_2x_3+2x_1x_3+3x_1x_4-x_3x_4=[x_1x_2+x_2x_3+2x_1x_3+3x_1x_4]-x_3x_4$\\
    $=(x_1+x_3)(x_2+2x_3+3x_4)-x_3(2x_3+3x_4)-x_3x_4$\\
    $=(1/4)[(x_1+x_3+x_2+2x_3+3x_4)^2-(x_1+x_3-(x_2+2x_3+3x_4))^2]-2x_3^2-4x_3x_4$\\
    $=(x_1+x_2+3x_3+3x_4 )^2/4-(x_1-x_2-x_3-3x_4 )^2/4-2(x_3+x_4 )^2+2x_4^2$.
}

\theorem{Proposition}{}{true}{
    Soit $E$ un $\mathbb{R}$-espace vectoriel de dimension finie de dimension $n$, et $q$ une forme quadratique sur $E$.\\
    Alors il existe $\ell_1, \ldots, \ell_p, \ldots, \ell_{p+r}$ des formes linéaires indépendantes telles que :
    \[
        q = \ell_1^2 + \cdots + \ell_p^2 - \ell_{p+1}^2 - \cdots - \ell_{p+r}^2
    \]
}
\remark{Cela découle directement de la décomposition de Gauss.}

\remark{
    Dans ce cas, $Mat_e(q) = Diag(\underbrace{1, \ldots, 1}_{p \text{ fois}}, \underbrace{-1, \ldots, -1}_{r \text{ fois}}, 0, \ldots, 0)$, et $rg(q) = rg(Mat_e(q)) = p + r$.
}

\subsubsection{Signature d'une forme quadratique}

\definition{
    Soit $q$ une forme quadratique sur $E$.\\
    On dit que $q$ est \textbf{positive} si $q(x) \geq 0$ pour tout $x \in E$, et que $q$ est \textbf{négative} si $q(x) \leq 0$ pour tout $x \in E$.\\
    On dit que $q$ est \textbf{définie positive} (resp. \textbf{définie négative}) si $q(x) > 0 \Leftrightarrow x \neq 0$ (resp. $q(x) < 0 \Leftrightarrow x \neq 0$).
}

\theorem{Théorème}{}{false}{
    Soit $E$ un $\mathbb{R}$-espace vectoriel de dimension finie $n$, et $q$ une forme quadratique sur $E$.\\
    Alors le couple $(p,r)$ d'entiers tels que $q = \ell_1^2 + \cdots + \ell_p^2 - \ell_{p+1}^2 - \cdots - \ell_{p+r}^2$ (pour $\ell_1, \ldots, \ell_{p+r}$ des formes linéaires indépendantes) ne dépend pas de $(\ell_1, \ldots, \ell_{p+r})$.
}

\theorem{Lemme}{}{false}{
    Soit $q$ une forme quadratique sur $E$ de dimension finie.\\
    Soit $(e_1, \ldots, e_n)$ une base q-orthogonale de $E$.\\
    \begin{enumerate}
        \item Si $q(e_i) \geq 0$ pour tout $i$, alors $q$ est positive,
        \item Si $q(e_i) > 0$ pour tout $i$, alors $q$ est définie positive.
    \end{enumerate}
}
\ndlr{cf. Laurent pour la preuve du lemme}
\ndlr{cf. Laurent pour la preuve du théorème}

\definition{
    Soit $E$ un $\mathbb{R}$-espace vectoriel de dimension finie, et $q$ une forme quadratique sur $E$.\\
    On appelle \textbf{signature} de $q$ le couple d'entiers $(p,r)$ tel que $q = \ell_1^2 + \cdots + \ell_p^2 - \ell_{p+1}^2 - \cdots - \ell_{p+r}^2$ pour $\ell_1, \ldots, \ell_{p+r}$ des formes linéaires indépendantes.\\

    De manière équivalente, la signature de $q$ est le couple d'entiers $(p,r)$ tel que pour toute base $q$-orthogonale $e$ de $E$, on a $p = \#\{ i \mid q(e_i) > 0 \}$ et $r = \#\{ i \mid q(e_i) < 0 \}$.
}

\example{
    \begin{itemize}
        \item Un produit scalaire est de signature $(dim(E), 0)$.
        \item $q(x,y,z,t) = x^2 + y^2 + z^2 - c^2 t^2$ est de signature $(3, 1)$.
        \item $q(x,y) = xy$ est de signature $(1, 1)$.
    \end{itemize}
}

\subsubsection{Liens entre signature et matrices symétriques}

Soit $q$ une forme quadratique sur $E$. Soit $\varphi$ sa forme polaire.\\
Soit $e=(e_1, \ldots, e_n)$ une base quelconque de $E$.\\
Posons $A = Mat_e(q) = Mat_e(\varphi) = (\varphi(e_i, e_j))_{1 \leq i,j \leq n}$.\\
Par définition $A$ est symétrique. Donc $A$ est représente la matrice d'un endomorphisme auto-adjoint.\\
Donc $\exists P \in O_n(\mathbb{R})$ tel que ${}^t P A P = Diag(\lambda_1, \ldots, \lambda_n)$, où $\lambda_1, \ldots, \lambda_n$ sont les valeurs propres de $A$.\\
Soit $(e_1', \ldots, e_n')$ la base de $E$ telle que $P = P_{e \to e'}$, \textit{i.e.} $e'_j = \sum_{i=1}^n P_{ij} e_i$ pour tout $j$.\\
Alors $e'$ est une base $q$-orthogonale de $E$, et $Mat_{e'}(q) = Diag(\lambda_1, \ldots, \lambda_n)$.\\
Donc, la signature de $q$ est donnée par le nombre de valeurs propres strictement positives et le nombre de valeurs propres strictement négatives de $A$.

\newpage
\tableofcontents

\end{document}